\subsection{Homotopy type in terms of critical values}

Let $f$ be a Morse function on a manifold $M$, then for $a\in\bR$ we define the {\emphcolor level set}
$$ M_a = f^{-1}(-\infty,a] = \set{p\in M}[f(p)\leq a] . $$
Because $(-\infty,a]$ has codimension $0$ in $\bR$, $f$ is trivially transervsal to it, and so $M_a$
is a smooth submanifold of $M$.
It turns out that these level sets offer rich insight into $M$.
For one, if $M$ is compact then $f$ attains a minimum $a$ and a maxmimum $b$, so that
$$ M_x = \emptyset\hbox{ when $x<a$},\qquad M_y = M\hbox{ when $y\geq b$} . $$
Now $a,b$ are critical values of $f$, and it turns out that these level sets have interesting behavior
around level sets.

For a Riemannian metric $g$ on $M$ and $f\in C^\infty(M)$, define $\grad f\in\X(M)$ to be the vector field
determined by
$$ \iota(\grad f)g = \d f . $$

\bthrm[name={nocritdiff}]

    Let $f$ be a smooth real-valued function on $M$, and let $a<b$ be real.
    Suppose that $f^{-1}[a,b]\subseteq M$ is compact and contains no critical points of $f$.
    Then $M_a$ and $M_b$ is diffeomorphic, and the inclusion $M_a\inj M_b$ is a strong deformation retract.

\ethrm

\bproof

    Choose a Riemannian metric on $M$, and define $\rho\colon M\to\bR$ to be a smooth function equal to
    $1/\iprod{\grad f,\grad f}$ on the compact set $f^{-1}[a,b]$ which vanishes outside a compact
    neighborhood of $f^{-1}[a,b]$.
    Note that $\rho$ is only-well defined by the assumption that $f$ has no critical points in $f^{-1}[a,b]$.
    Then define $X=\rho\grad(f)$.
    This vector field is compactly supported and thus generates a global flow $\theta_t\colon M\to M$.

    Fix a $p\in M$ and consider the map $t\mapsto f(\theta_t(p))$, which is smooth $\bR\to\bR$.
    Then
    $$ \drvof{f(\theta_t(p))}t = \iprod{\drvof{\theta_t(p)}t,\grad f} = \iprod{X,\grad f} =
    \frac\rho{\grad f,\grad f} . $$
    This is $1$ when $a\leq f(\theta_t(p))\leq b$, and otherwise we may assume it is bound by $1$.
    Thus $t\mapsto f(\theta_t(p))$ is linear of the form $t+f(p)$ when $a\leq f(\theta_t(p))\leq b$, and
    otherwise is less than $t+f(p)$.

    Consider the diffeomorphism $\theta_{b-a}\colon M\to M$.
    Then if $f(p)\leq a$ we have $f(\theta_{b-a}(p))\leq b-a+f(p)\leq b$ so $\theta_{b-a}(M_a)\subseteq M_b$.
    Similarly we see that $\theta_{a-b}=\theta_{b-a}^{-1}$ maps $M_b$ into $M_a$ so that they form a
    diffeomorphism $M_a\cong M_b$.

    Now define a homotopy $H\colon[0,1]\times M_b\to M_b$ by
    $$ H_t(p) = \cases{p&if $f(p)\leq a$\cr \phi_{t(a-f(p))}(p)&if $a\leq f(p)\leq b$} . $$
    Both of these cases agree when $f(p)=a$, and thus this is continuous.
    Now $H_0$ is the identity, $H_t(p)=p$ when $p\in M_a$, and
    $H_1(p)=\phi_{a-f(p)}(p)\in M_a$ for $p\notin M_a$ so that $H$ is a strong deformation retract.
    \qed

\eproof

\bthrm

    Let $f\colon M\to\bR$ be a smooth function and $p\in M$ a nondegenerate critical point with index
    $\lambda$ and $f(p)=c$.
    If $f^{-1}[c-\epsilon,c+\epsilon]$ is compact and contains no critical points other than $p$ then
    $\epsilon$ can be taken sufficiently small so that $M_{c+\epsilon}$ is homotopy equivalent to
    $M_{c-\epsilon}$ attached to a $\lambda$-cell.

\ethrm

\bnote

    By a $\lambda$-cell we mean a manifold diffeomorphic to the unit ball in $\bR^\lambda$.
    Attaching it to a manifold means gluing it along some $g\colon\partial e_\lambda\to M$.

\enote

\bproof

    Take a neighborhood $U$ around $p$ such that $f$ has the representation
    $$ f = c - u_1^2 - \cdots - u_\lambda^2 + u_{\lambda+1}^2 + \cdots u_n^2 . $$
    Take $\epsilon>0$ sufficiently small so that $f^{-1}[c-\epsilon,c+\epsilon]$ contains no critical points
    other than $p$ and $U$ contains the closed ball $\set{\bar u}[\norm u^2\leq2\epsilon]$.
    Define $e_\lambda$ in $U$ to be
    $$ e_\lambda = \set{(u_1,\dots,u_\lambda,0,\dots,0)\in U}[u_1^2+\cdots+u_\lambda^2\leq\epsilon] . $$
    Now, $\bar u\in e_\lambda\cap M_{c-\epsilon}$ iff
    $$ u_{\lambda+1} = \cdots = u_n = 0,\qquad u_1^2+\cdots+u_\lambda^2\leq\epsilon,\qquad
    -u_1^2+\cdots+u_\lambda^2 \leq -\epsilon , $$
    i.e.\ $e_\lambda\cap M_{c-\epsilon}=\partial e_\lambda$.
    Thus in $U$ we see that $e_\lambda\cup M_{c-\epsilon}$ is a gluing of $e_\lambda$ and $M_{c-\epsilon}$.

    Let $\mu\colon\bR\to\bR$ be a smooth function such that
    $$ \mu(0) > \epsilon,\qquad \mu(r) = 0 \hbox{ for $r\geq2\epsilon$},\qquad -1 < \mu' \leq 0 . $$
    Define $F\colon M\to\bR$ to be equal to $f$ outside of $U$ and within $U$,
    $$ F = f - \mu(u_1^2 + \cdots + u_\lambda^2 + 2u_{\lambda+1}^2 + \cdots + 2u_n^2) . $$
    This is a smooth function as it agrees on its overlaps.
    For convenience define
    $$ \xi,\eta\colon U\to[0,\infty),\qquad \xi = u_1^2 + \cdots + u_\lambda^2,\qquad
    \eta = u_{\lambda+1}^2 + \cdots + u_n^2 $$
    so that
    $$ f = c - \xi + \eta,\qquad F = c - \xi + \eta - \mu(\xi + 2\eta) . $$

    Now we make a few claims which will progress our proof.
    \benum
        \item $F^{-1}(-\infty,c+\epsilon]=M_{c+\epsilon}=f^{-1}(-\infty,c+\epsilon]$.
        Outside the ellipsoid $\xi+2\eta\leq2\epsilon$ both $f$ and $F$ agree, so that
        these sets agree outside of this ellipsoid.
        Within the ellipsoid note that
        $$ F \leq f = c - \xi + \eta \leq c + \frac12\xi + \eta \leq c + \epsilon $$
        and so the ellipsoid is contained in both sets.

        \item The critical points of $F$ and $f$ are the same.
        We have
        $$ \pdvof F\xi = - 1 - \mu'(\xi + 2\eta) < 0,\qquad
        \pdvof F\eta = 1 - 2\mu'(\xi + 2\eta) \geq 1 , $$
        and so
        $$ \d F = \pdvof F\xi\d\xi + \pdvof F\eta\d\eta , $$
        and since $\d\xi,\d\eta$ are only every simultaneously zero at the origin, $F$'s only critical
        point in $U$ is the origin.
        And outside of $U$ it agrees with $f$.

        \item $F^{-1}(-\infty,c-\epsilon]$ is a deformation retract of $M_{c+\epsilon}$.
        Now consider $F^{-1}[c-\epsilon,c+\epsilon]=F^{-1}[c-\epsilon,\infty)\cap F^{-1}(-\infty,c+\epsilon]$.
        By the first point, the second set here is $f^{-1}(-\infty,c+\epsilon]$.
        And since $F\leq f$ in $U$, the first set is contained in $f^{-1}[c-\epsilon,\infty)$.
        Putting this together we have
        $$ F^{-1}[c-\epsilon,c+\epsilon] \subseteq f^{-1}[c-\epsilon,c+\epsilon] . $$
        Thus it is compact, and by (2) the only critical point of $F$ in it is potentially the origin.
        But
        $$ F(p) = c - \mu(0) < c - \epsilon , $$
        so that the origin is not in $F^{-1}[c-\epsilon,c+\epsilon]$.
        By \refmath[theorem]{nocritdiff} this means that $F^{-1}(-\infty,c-\epsilon]$ is a deformation
        retract of $F^{-1}(-\infty,c+\epsilon]$ which is $M_{c+\epsilon}$ by (1).
    \eenum
    Let us denote by $H$ the closure of $F^{-1}(-\infty,c-\epsilon]-M_{c-\epsilon}$ so that
    $F^{-1}(-\infty,c-\epsilon]=M_{c-\epsilon}\cup H$.
    Now we have defined
    $$ e_\lambda = \set{q\in U}[\xi(q)\leq\epsilon,\ \eta(q)=0] . $$
    Because $\pdvof F\xi<0$ we have that for $q\in e_\lambda$, $F(q)\leq F(p)<c-\epsilon$
    and $f(q)\geq c-\epsilon$.
    Thus $e_\lambda$ is contained in $F^{-1}(-\infty,c-\epsilon]-M_{c-\epsilon}\subseteq H$.

    Finally we claim that $M_{c-\epsilon}\cup e_\lambda$ is a deformation retract of $M_{c-\epsilon}\cup H$.
    In lieu of (3) this proves the desired result.
    The deformation retract $h_t$ is defined as follows:
    \benum
        \item For $\xi\leq\epsilon$ define
        $$ h_t(u_1,\dots,u_n) = (u_1,\dots,u_\lambda,tu_{\lambda+1},\cdots,tu_n) $$
        so that $h_0$ maps into $e_\lambda$, $h_1$ is the identity, and $h_t|_{e_\lambda}$ is the identity.
        This is well-defined, as in it maps $F^{-1}(-\infty,c-\epsilon]$ to itself, because
        $\pdvof F\eta>0$ so that $F(\xi,t\eta)\leq F(\xi,\eta)\leq c-\epsilon$ for $t\leq1$.
        \item For $\epsilon\leq\xi\leq\eta+\epsilon$, let
        $$ h_t(u_1,\dots,u_n) = (u_1,\dots,u_\lambda,s_tu_{\lambda+1},\dots,s_tu_n),\qquad
        s_t = t + (1-t)\parens{\frac{\xi-\epsilon}\eta}^{1/2} . $$
        Here $s_t\in[0,1]$ so that this is well-defined for the same reason as above.
        And $s_1=1$ so that $h_1$ is the identity, and $h_0$ maps into
        $f^{-1}(c-\epsilon)\subseteq M_{c-\epsilon}$.
        When $\xi=\epsilon$ this agrees with the definition in (1), so it remains continuous.
        \item For $\eta+\epsilon\leq\xi$, i.e.\ within $M_{c-\epsilon}$, define $h_t$ to be the identity.
        When $\xi=\eta+\epsilon$ this corresponds to the definition in (2).
    \eenum
    And so we have finished.
    \qed

\eproof

The following corollary can be obtained from topological considerations regarding, and is immediate in the
case when $M$ is compact (by induction).

\bcoro

    Let $f\colon M\to\bR$ be a Morse function on $M$ such that for every $a\in\bR$, $M_a$ is compact.
    Then $M$ is homotopy-equivalent to a CW-complex with a cell of dimension $\lambda$ for each critical
    point of index $\lambda$.

\ecoro

